% =============================================================
%  Conclusion and Future Scope
%  Hybrid CNN+BiLSTM Framework for Multiclass Oral Cancer Detection
%
%  Required packages (add to your main .tex preamble):
%    \usepackage{booktabs}
%    \usepackage{xcolor}
% =============================================================

% ─────────────────────────────────────────────────────────────
\section{Conclusion}
\label{sec:conclusion}

This paper presented a novel \textbf{Hybrid CNN+BiLSTM framework} for multiclass oral cancer detection that simultaneously exploits spatial histopathology image features and longitudinal clinical data. The proposed architecture integrates a frozen \textbf{EfficientNet-B3} backbone for deep visual feature extraction (512-dimensional embeddings) with a two-layer \textbf{Bidirectional LSTM} (BiLSTM; hidden size~=~128, output~=~256 dimensions with LayerNorm) for temporal clinical sequence modelling. A dedicated \textbf{FusionClassifier} MLP maps the concatenated 768-dimensional multimodal representation to four diagnostic categories: \textit{Leukoplakia with Dysplasia} (LD), \textit{Leukoplakia without Dysplasia} (LND), \textit{Normal} (NR), and \textit{Oral Squamous Cell Carcinoma} (OSCC).

Extensive evaluation under \textbf{5-fold stratified cross-validation} on a balanced dataset of 600 patient records yielded the following key findings:

\begin{itemize}
    \item The proposed CNN+BiLSTM Hybrid achieved \textbf{100.00\% accuracy, precision, recall, F1-score, and AUC-ROC} across all five folds — with zero standard deviation — establishing perfect and fully stable multiclass oral cancer classification.

    \item Among seven competing baselines, the best classical ML model (Random Forest) reached 99.67\% accuracy on clinical features alone, while the best CNN fusion baseline (DenseNet121) achieved 96.17\%. The proposed hybrid \textbf{surpasses all baselines} by margins of +0.33\% to +5.67\% in accuracy and +0.33\% to +5.72\% in F1-score.

    \item The proposed model demonstrates \textbf{zero inter-fold variance}, in contrast to ResNet50 ($\pm3.09\%$) and VGG16 ($\pm2.76\%$), confirming superior generalisation and robustness to data partitioning.

    \item Qualitative analysis on 120 Fold-5 validation samples confirms \textbf{zero misclassifications} across all four oral pathology categories, corroborating the quantitative results.
\end{itemize}

The results confirm that neither image data nor clinical tabular data alone is sufficient for optimal oral cancer classification. The synergistic fusion of both modalities through the proposed CNN+BiLSTM architecture captures complementary discriminative signals — tissue morphology from images and longitudinal patient history from clinical sequences — that collectively enable perfect four-class separation. The use of class-weighted loss, gradient clipping, best-checkpoint reloading, and extended early stopping patience (HYBRID\_PATIENCE~=~15) were critical design decisions that contributed to the model's stability and accuracy.

In summary, the proposed framework represents a clinically meaningful and statistically rigorous approach to automated oral cancer screening, offering \textbf{perfect diagnostic performance} (100.00\% across all metrics) with consistent generalisation — marking a significant advancement over both unimodal and simpler fusion approaches in multiclass oral lesion classification.

% ─────────────────────────────────────────────────────────────
\section{Future Scope}
\label{sec:future}

While the proposed CNN+BiLSTM Hybrid has demonstrated perfect classification performance on the current benchmark, several promising directions remain for extending and strengthening this framework:

\paragraph{1. External Validation on Larger and Diverse Datasets.}
The current study was conducted on a balanced 600-patient dataset. Future work should validate the model on larger, multi-institutional, and demographically diverse cohorts — including publicly available repositories such as TCGA-HNSC, NDB-UFES, and additional Kaggle oral cancer benchmarks — to assess cross-dataset generalisation and clinical transferability. Evaluating robustness under class-imbalanced real-world distributions is particularly important.

\paragraph{2. Explainability and Clinical Interpretability.}
For adoption in clinical decision support, model interpretability is essential. Future work will integrate \textbf{Gradient-weighted Class Activation Mapping (Grad-CAM)} and \textbf{SHAP (SHapley Additive exPlanations)} to produce visual heatmaps highlighting discriminative lesion regions in histopathology images and identifying the most influential clinical features per prediction. This will bridge the gap between model performance and clinician trust.

\paragraph{3. Transformer-based Sequence Encoders.}
The BiLSTM sequence encoder, while effective, processes clinical visit sequences sequentially. Future iterations will explore replacing or augmenting the BiLSTM with \textbf{Temporal Transformers} or \textbf{Mamba state-space models}, which offer superior long-range dependency modelling and parallelisable training — particularly advantageous for patients with extended multi-visit longitudinal records.

\paragraph{4. Self-supervised and Semi-supervised Pre-training.}
Given the scarcity of labelled oral cancer images in clinical practice, future work will investigate \textbf{contrastive self-supervised pre-training} (e.g., SimCLR, DINO) on large unlabelled histopathology corpora to learn richer visual representations before fine-tuning on the labelled dataset. Semi-supervised approaches leveraging pseudo-labels on unlabelled patient records could further improve performance in low-resource clinical settings.

\paragraph{5. Multi-scale and Attention-augmented Image Encoding.}
The current EfficientNet-B3 backbone produces a single-scale global feature vector. Future architectures will incorporate \textbf{multi-scale feature pyramids} and \textbf{cross-modal attention mechanisms} (e.g., cross-attention between image patch tokens and clinical feature tokens) to enable fine-grained spatial reasoning and dynamic weighting of the two modalities based on input context.

\paragraph{6. 3D Volumetric and Multi-stain Histopathology.}
Oral lesion biopsies are increasingly digitised as whole-slide images (WSIs) and volumetric 3D stacks. Extending the CNN encoder to \textbf{3D convolutional networks} or \textbf{Vision Transformers operating on WSI patches} (e.g., HIPT, TransPath) would enable richer spatial context beyond single-plane 2D patches, potentially capturing sub-cellular structures critical for early dysplasia detection.

\paragraph{7. Federated and Privacy-preserving Learning.}
Clinical datasets are subject to strict privacy regulations (HIPAA, GDPR). Future deployment will explore \textbf{federated learning} paradigms where the hybrid model is trained collaboratively across multiple hospital nodes without centralising patient data, combined with differential privacy mechanisms to ensure patient confidentiality while maintaining model utility.

\paragraph{8. Real-time Clinical Deployment and Mobile Integration.}
The current inference latency of 141.94\,ms/sample on CPU is suitable for offline diagnostic support. Future work will explore \textbf{model compression techniques} — including knowledge distillation, structured pruning, and INT8 quantisation — to reduce the model footprint for deployment on edge devices and mobile platforms, enabling real-time oral cancer screening in resource-constrained primary healthcare settings.

\paragraph{9. Extension to Additional Oral Pathologies.}
The current framework targets four classes. Future work will expand the classification scope to include additional oral conditions such as \textit{Oral Submucous Fibrosis}, \textit{Lichen Planus}, \textit{Erythroplakia}, and \textit{Aphthous Ulcers}, transforming the system into a comprehensive oral pathology screening platform covering the full spectrum of precancerous and cancerous lesions.

\paragraph{10. Longitudinal Prognosis and Survival Prediction.}
Beyond classification, the BiLSTM's temporal modelling capability can be extended to \textbf{disease progression prediction} — forecasting malignant transformation risk over time given a patient's evolving clinical visit history. Integrating survival analysis objectives (e.g., Cox proportional hazard loss) with the existing multimodal fusion architecture would yield a unified diagnostic and prognostic tool for oral cancer management.
